\section{Research frontiers and open problems}\label{sec:frontiers}


\regime{Control}{the open problems on the return arc, and concrete programs to close them.}
The preceding parts establish a diagnosis, not a celebration. Across the
$N{=}435$ coded corpus, the works we examined have learned how to accumulate state
and how to retrieve it: substrates, write pipelines, retrieval rankers,
consolidation schedules, and skill libraries are well studied and heavily benchmarked. They
are less developed on how to \emph{govern} state once it persists. The two thinnest
cells in the coding scheme are the governance axes most tied to consequential
state: only $27/435$ works expose any rollback mechanism, and authority is
the rarest of the six state axes at $72/435$. Every part of this survey reported
the same pattern under a different name. P1 finds almost no foundational work on formal
authority structure and zero on state recovery after corruption; P3 has zero works
that directly address rollback and only two that address audit lineage; P5 states
its deletion and rollback semantics abstractly because the mechanistic literature to
ground them does not yet exist; P6 lacks formal consistency theory for shared
mutable state; P7 lists almost all of its weak cells as ``no benchmark exists'' for
the very governance properties the thesis prioritizes. The agenda is broader than
 retrieval accuracy alone: it is the work of converting an
accumulate-and-retrieve framing into a govern-over-a-lifecycle framing.

One observation should temper how this frontier is read. The governance mechanisms
the survey was able to cite, origin-bound write authority, semantic transactions,
saga-style compensation, audited skill graphs, runtime authority gating, are
overwhelmingly 2025 and 2026 preprints, and few have appeared in widely deployed,
long-running systems or been measured under realistic fault load. That clustering
admits two readings, and the honest position holds both. The optimistic reading is
that the field is correctly identifying the gap and mobilizing quickly, so the
preprint surge is genuine frontier work on a normal multi-year maturation timeline.
 The cautious reading is that governance may be harder to build and benchmark than
memory: it demands formal reasoning, cross-system testing, and infrastructure that a
paper can propose but not easily scale, so some of these mechanisms may not survive
contact with production. We report the frontier as we find it, and treat its youth as
itself a signal: the direction is right, but durability is unproven, and the agenda
is therefore stated as work to be done rather than results to be consolidated.

Five coupled research programs follow from that diagnosis. The first four are the
directions the coverage map already implies, each anchored to a quantified corpus
gap and stated as a coupled pair of a measurement instrument and a mechanism the
instrument would discipline. The fifth is the meta-observation that makes the
other four tractable: late rounds of in-vocabulary ``agent memory'' queries showed
diminishing returns in our search frame, so the governance backbone must borrow
native vocabulary, and in many cases entire correctness calculi, from databases
and distributed systems, formal methods, capability security, machine unlearning,
and regulatory compliance. Each direction below gives the gap, why it matters for
persistent state, a concrete program of work, and the open problem it leaves
standing. Table~\ref{tab:frontier-map} ties the five programs to the corpus cells
they target and to the closest current anchors.

\begin{table}[t]
\centering
\small
\caption{The five frontier programs, the quantified corpus gap each targets, the
closest current anchors in the corpus, and the open instrument each program needs.
The recurring pattern is that mechanisms exist as point solutions while the
measurement and the cross-axis composition are absent.}
\label{tab:frontier-map}
\begin{tabular}{p{2.3cm}p{2.4cm}p{2.6cm}p{2.4cm}}
\toprule
\textbf{Program} & \textbf{Corpus gap} & \textbf{Closest anchors} & \textbf{Missing instrument} \\
\midrule
Lifecycle-complete evaluation &
Rollback $27/435$; recovery never measured &
\citet{chang2025sagallm}, \citet{zheng2026acrfence} &
Recovery-success and recovery-cost metrics on hosted systems \\
\addlinespace
Controlled compounding &
Consolidation drops retrieval handles; contagion has no safe threshold &
\citet{kang2026oslmr}, \citet{liu2026memorycontagion} &
Write-acceptance gate keyed to later addressability \\
\addlinespace
Authority, privacy, deletion, cross-surface &
Authority $72/435$; no single-policy cross-surface benchmark &
\citet{lahjouji2026privacysurvey}, \citet{shah2026unlearningmirage} &
One agent, all surfaces, one declared policy \\
\addlinespace
Shared-memory governance &
No distributed rollback; revocation cascades unmodeled &
\citet{ren2026gatemem}, \citet{zhu2026openport} &
Authority-scoped rollback of tool-originated state \\
\addlinespace
Adjacent-discipline bridging &
Governance axes thin in every part &
\citet{park2026kumiho}, \citet{parakhin2026bureaucracy} &
Native correctness calculi imported, not re-derived \\
\bottomrule
\end{tabular}
\end{table}

\subsection{Lifecycle-complete evaluation and running AOEP on real systems}

The highest-risk gap in the corpus is not that rollback mechanisms are
rare but that rollback is studied as a \emph{mechanism} and almost never as a
\emph{measured capability}. The $27/435$ figure undercounts the problem: even among
those twenty-seven works, none reports the quantities an operator of an always-on
agent would actually need, namely how often recovery succeeds after a corruption
event, how much state is lost when it does, and what the recovery costs in latency
and compute. The vocabulary for these quantities already exists in databases and
distributed systems as recovery-point and recovery-time objectives, mean-time-to-recovery,
and recovery-success rate, but the agent-memory literature has not adopted it. The
result is a field that can describe checkpoint-and-restore designs while being unable
to say whether any of them works under realistic fault load.

The mechanism literature that does exist is suggestive but partial. \citet{chang2025sagallm}
imports the database Saga model into multi-agent LLM planning, giving compensable
execution, automated compensation, and independent validation agents so that a
multi-step workflow can be rolled back through compensating actions rather than naive
undo. This is the right import, but it assumes compensation handlers exist and that
the prior state can be reconstructed, the assumption that fails for irreversible
side effects. \citet{zheng2026acrfence} confronts that failure directly:
it characterizes semantic rollback attacks in agent checkpoint-restore, where naive
restore either duplicates an already-consumed effect or resurrects an authority that
was meant to lapse, and mitigates by recording irreversible effects and choosing
replay-or-fork accordingly. \citet{nakajima2026regimes} offers a third stance: an
event-sourced runtime where state is a deterministic projection of an append-only
log, so exact replay is free, at the cost of an ever-growing log with no garbage
collection. \citet{shawn2026pace} attacks the commit decision rather than the
restore, replacing greedy score-based acceptance of self-updates, which p-hacks
itself into harmful commits, with an anytime-valid betting gate that bounds the
probability of a false commit. These four works disagree about where recovery should
live: in compensation logic, in effect ledgers, in the log substrate, or in the
commit gate. None of them is evaluated against the others on a common recovery
benchmark in the corpus, because we found no shared benchmark for this case.

The concrete program is to make recovery a first-class measured property and to run
the measurement on real systems rather than on synthetic stores. An always-on
evaluation harness should inject corruption, poisoned writes, stale facts, partial
write failures, and lapsed-authority retries, then score recovery-success rate,
state lost at the recovery point, and recovery cost, against a stated objective. The
critical design constraint, drawn from \citet{zheng2026acrfence}, is that the harness
must distinguish reversible from irreversible state, because a restore that looks
correct on retrievable memory can still have duplicated an external side effect. The
open problem this leaves is sharp: within our query frame, there is no agreed definition of a
``recovered'' always-on agent. \citet{orogat2026gem} formalizes governed evolving
memory with six correctness conditions, including rollback traceability, and proves
that record-level stores cannot satisfy them, but it is a theoretical framework
without an empirical harness, and it does not cover the authority and provenance
policies that determine whether a given restore is even permitted. Until recovery is
measured on hosted, multi-task agent benchmarks rather than asserted in design
sections, the thesis's recoverability axis remains a property the current
literature often claims but rarely demonstrates.

\subsection{Controlled compounding and retrieval-addressable consolidation}

Always-on agents are supposed to get better with use. Replay, reflection, skill
libraries, and consolidation are the mechanisms by which experience is meant to
compound into competence. The plasticity-stability tension is well studied as a
learning problem, but as a \emph{governance} problem it exposes a different failure:
the same accumulation that should compound verified competence also lets bad
trajectories, stale facts, and poisoned traces compound into durable policy. The
question is not how to remember more but how to admit only what should become
permanent, and to keep what is admitted addressable by the queries that will later
need it.

Two distinct hazards live at the consolidation boundary. The first is loss of
addressability. Consolidation compresses, and compression can score well on aggregate
recall while silently dropping the retrieval handles, the names, identifiers, dates,
and keys, that a later query must match; a memory then looks compact and coherent
while becoming unusable. The principle the agenda names for this is
retrieval-addressable consolidation: accept a write or update only when it preserves
the handles a later query will use, rather than only preserving topical gist. The
second hazard is contamination that compounds across time. \citet{liu2026memorycontagion}
shows that evaluator bias embedded in stored trajectories propagates
cross-temporally to future agents that share the memory, and that there is no safe
contamination threshold even under oracle consolidation. This is a strong warning
against naive compounding: without an admission gate, shared experience can be
systematically harmful rather than only occasionally noisy. The retention literature gives the lever such a gate would pull.
\citet{kang2026oslmr} casts retention as an NP-hard constrained stochastic
optimization over budget, usefulness, and delayed miss-or-stale cost, learning a
query-conditioned, observability-safe policy that beats recency and heuristic
baselines under tight budgets, the regime where addressability is
sacrificed first. \citet{kumar2026memarchitect} adds a rule-based layer that governs
decay, resolves contradictions among facts, and filters stale ``zombie'' memories
before they re-enter retrieval.

What none of these works supplies is a write-acceptance gate that jointly enforces
both addressability and provenance-bounded admission as a precondition for a write
becoming durable. \citet{shawn2026pace} provides the statistical shape of such a
gate for self-updates, an anytime-valid betting bound on false commits, but does not
key it to whether the committed item remains retrievable, and \citet{joshi2026eywa}
provides the provenance shape, storing immutable source evidence first and deriving
canonical facts only after validation against typed signals, but does not measure
addressability after derivation. The open problem is to unify them: a consolidation
operator that compounds competence only when the candidate write is both
authorized-by-provenance and addressable-by-future-query, and that is evaluated on
whether long-horizon competence actually rises rather than on single-step recall.
The contagion result of \citet{liu2026memorycontagion} makes the stakes explicit.
A field that consolidates for recall without an admission gate is building durable
policy out of unvetted experience, which is the plasticity-stability invariant
failing in its most damaging direction.

\subsection{Authority, privacy, deletion, and a cross-surface benchmark}

The authority axis is the rarest in the corpus at $72/435$, and its scarcity is not
an accident of coding: authority, the question of who may read, write, delete, or
act on a piece of state, is the axis that the accumulate-and-retrieve paradigm has no
reason to represent. An always-on agent that never forgets and never checks who
licensed a fact will treat a preference that was true last month, changed yesterday,
as still authoritative today, and will let any stored fact influence any action. The
governance program here has three tightly linked obligations, privacy, deletion, and
authority, and a missing instrument that would test them together.

On privacy, the corpus has moved from training-time leakage to deployment-time and
memory-resident leakage. \citet{carlini2021extractingtrainingdata} established that
models memorize and leak verbatim training data; \citet{chen2026deploymenttimememorization}
shows the analogous failure for deployed agents memorizing sensitive deployment-time
data, and \citet{chen2026mrmmia} gives the first membership-inference framework aimed
explicitly at chat-agent persistent memory units. \citet{naseh2025riddle} shows the
attack needs only about thirty natural-language queries to infer datastore
membership while evading query-rewriting defenses, and \citet{xu2026memorysilent}
isolates the deeper failure with its RBI-Eval probe: agents integrate sensitive
memory content even when it is not warranted, and retrieval cannot prevent
integration once the memory reaches the generator. On deletion, the picture is
worse, because deletion is widely assumed to be solved and is not.
\citet{staufer2025whatshouldllmsforget} addresses the prerequisite everyone else
assumes away, identifying \emph{what} must be forgotten under right-to-be-forgotten
obligations; \citet{xue2025unlearningverificationsurvey} surveys how one would
\emph{verify} that deletion succeeded, mapping directly onto the recoverability axis;
and \citet{shah2026unlearningmirage} gives a negative result, showing
that supposedly forgotten information resurfaces under multi-hop and aliased queries.
\citet{wang2024unlearningmeetsrag} and \citet{koga2024dprag} give the two
mechanism families, unlearning over retrieval stores and differential privacy spent
only on sensitive tokens, but neither closes the cascade: deleting a source fact does
not delete the summaries, embeddings, and derived facts it spawned.
\citet{zhao2026memorepair} is the closest corpus work to a cascade solution, a
barrier-first repair contract that withdraws stale derived descendants when a source
is deleted and republishes only validated, predecessor-closed successors, cutting
invalidated-memory exposure toward zero; it should be read as the deletion-propagation
invariant made operational, and as evidence of how rare such operationalization
currently is.

On authority, several recent works treat it as a property that must only
narrow. \citet{zhu2026intentgovernedauth} makes the user's expressed intent a
monotone, session-scoped policy that can only reduce the authority static credentials
grant; \citet{ibrahim2026overlaygovernance} proves how delegated authority attenuates
along recursive delegation chains; and \citet{parakhin2026bureaucracy} establishes a
structural equivalence between hardware memory-consistency models and authorization
revocation, bounding unauthorized operations by execution count rather than by a
time-to-live that an idle agent can outlast. \citet{dash2026selfportrait} supplies
the empirical motivation from the user's side: an audit of 2{,}050 real ChatGPT
persistent-memory entries found 96\% were silently system-created, so the authority
over what an agent remembers about a user is, in practice, held by the agent and not
the user. \citet{zhang2025ragmemoryperceptions} and \citet{feng2025regulatoryui}
confirm from interviews and system analysis that users hold incomplete mental models
of persistent memory and want granular rights to review, edit, and delete, with the
interface itself acting as a governance surface.

The missing instrument unifies all three obligations. The corpus contains no
benchmark that drives a single agent across all the data surfaces it touches, the
retrieval store, the tool APIs, the long-term memory, and the inter-agent channel,
under one declared privacy policy, and scores whether the policy holds end to end.
\citet{lahjouji2026privacysurvey} identifies this absence, finding that
information-flow control is the only governance that covers cross-session inference
leakage and that no existing benchmark exercises an agent across surfaces under one
policy. This is the cross-surface benchmark the agenda calls for: a privacy and
authority policy declared once, an agent that must honor inspection, correction,
erasure, and rollback over masked, access-controlled, and origin-bound memory, and a
score that fails the moment a deleted fact resurfaces in any surface or an action
descends from an authority that has lapsed. The open problem is that until such an
instrument exists, the deletion and authority literatures remain a collection of
surface-local point defenses, and the thesis's scope-non-expansion and
deletion-propagation invariants cannot be tested in the multi-surface setting where
they actually fail.

\subsection{Shared-memory governance and authority-scoped rollback}

Shared and multi-agent state multiplies every governance hazard, because a single
write can now influence principals who never authored it, and because the authority
that licensed a write may belong to one agent while the consequences land on another.
Two corpus cells stay persistently thin here. The first is the intersection of
authority and forgetting: who may compel deletion from shared memory, and who
verifies that it propagated to every principal's view. The second, and sharper, is
distributed rollback of tool-originated state once the authority that licensed the
original action has lapsed.

On the access-control side the corpus is maturing. \citet{ren2026gatemem}
benchmarks whether an LLM agent governing a shared memory pool can serve legitimate
requests while enforcing per-role and per-scope access controls and reliably
forgetting data on request across four real-world domains, which is the closest
existing instrument to a shared-memory governance benchmark. \citet{masoor2025samep}
and \citet{ravindran2026portablemem} build the cross-session and cross-vendor
sharing substrates, with AES-256-GCM access control and capability-scoped,
injection-resistant rehydration respectively, but each new shared channel widens
the authority and provenance surface that attacks exploit.
\citet{dalugoda2026hdp} contributes an append-only signed delegation-provenance chain
that binds each agent action to an originating human authorization and scope with
offline audit verification, which is the lineage half of authority-scoped rollback.
\citet{wright2025immutablemem} pushes immutability to its logical end, an append-only
Merkle ledger with privilege-lattice access control, but in doing so it collides head
on with deletion: a strictly immutable ledger cannot honor a forget request, so
immutability and deletion-propagation are in direct tension, and the works we
examined do not yet resolve which yields.

The authority-scoped rollback problem is where the corpus is thinnest. Consider the
canonical case: an agent makes a tool call under
a delegated authority, the side effect is irreversible or externally visible, the
authority later lapses or is revoked, and the call is retried after a transient
failure. A correct system must not let the retry resurrect already-consumed authority
or duplicate the irreversible effect. \citet{zhu2026openport} addresses the temporal
half of this with a State Witness profile that revalidates execution-time
preconditions and fails closed on state mismatch, killing stale approvals when
world-state drifts between an action's approval and its delayed execution, the
time-of-check-to-time-of-use drift that always-on agents create by design.
\citet{deochake2026heartbeatcred} addresses the liveness half: descendant agent
authority self-expires within a provably bounded window once parent liveness ceases,
killing zombie agents that keep acting after operator shutdown.
\citet{bhatt2025etdi} and \citet{south2025authenticateddelegation} bind grants to
signed tool definitions and to scoped, accountable delegation chains, so a silently
redefined tool or an over-broad delegation can be detected. But none of these works,
and no benchmark in the corpus, scores the full authority-scoped rollback case: a
retried call whose licensing authority has lapsed, against a shared substrate, with
an irreversible effect, where the correct outcome is neither blind replay nor blind
restore. This is the multi-agent instantiation of the recoverability axis, and with
only $27/435$ works exposing any rollback at all, it is unmeasured. Trigger
calibration, deciding when an always-on agent should interrupt or act on shared
state proactively rather than wait, is the related open axis: surfacing state too
eagerly is itself an authority overreach, and the utility-calibrated timing of
proactive action remains largely unstudied as a governance question rather than a
helpfulness one.

\subsection{Bridging from adjacent disciplines}

The four programs above share a constraint that the gap analysis makes visible:
late in-vocabulary agent-memory searches returned diminishing new governance
coverage under our query frame. The governance backbone is therefore unlikely to be
filled by re-running the same vocabulary. It needs deliberate import of native
correctness calculi from five established disciplines, treating their decades of
accumulated theory as part of the foundation the governance axes need. This bridging
is already visible in several strong recent works, which is evidence that the move
is productive.

From databases and distributed systems comes transactional correctness and recovery.
\citet{chang2025sagallm} imports the Saga model wholesale; the further frontier is
snapshot isolation, write-ahead logging, point-in-time recovery, and consistency
models, sequential, causal, and eventual, applied to concurrent writes on shared
agent memory, where P6 explicitly lacks formal consistency theory. From classical
knowledge representation comes belief revision: \citet{park2026kumiho} builds a
graph-native versioned memory whose update and contraction operations are proven to
satisfy the AGM and Hansson belief-base postulates, importing rationality guarantees
that the heuristic consolidation literature lacks entirely, and
\citet{wang2026toki} recasts contradiction resolution as write-time concurrency
control with a bitemporal operator algebra and four soundness theorems, giving the
both-truths representation an always-on agent needs to mark a fact superseded rather
than delete it. From formal methods and runtime verification comes provable
enforcement: \citet{bollig2026causalpast} defines a past-time temporal logic with a
provably correct vector-clock monitor for guards over causally-visible stored state;
\citet{winston2026solveraided} compiles natural-language tool-use policies into
SMT-LIB and uses Z3 as a per-call precondition gate; and
\citet{metere2026skillcontainment} promotes a skill manifest to a mechanically
checkable capability-containment proof. These works move enforcement from the design
section into the forward pass, which P5 flagged as its sparsest cell.

From capability security comes the discipline of authority that only narrows.
\citet{ibrahim2026overlaygovernance} overlays recursive delegation, dynamic scoping,
and resource-scope attenuation onto existing relational authorization with formal
proofs of attenuation along delegation chains; \citet{jin2026capseal} keeps raw
credentials non-exportable behind schema-checked invocations to defeat prompt-injection
exfiltration; and \citet{metere2026attestedtoolserver} admits tool servers only via
offline-signed clearance against a pinned trust root. These are capability-system
primitives, object capabilities, attenuation, confinement, transplanted into the
agent setting, and they are why authority, though rarest in the corpus, has
a plausible path toward stronger coverage. From machine unlearning and compliance comes the deletion
and audit backbone: \citet{nguyen2024machineunlearningsurvey} and
\citet{xue2025unlearningverificationsurvey} supply the methods and the verification
theory, \citet{shah2026unlearningmirage} supplies the adversarial test that any
import must pass, and the regulatory framings of \citet{staufer2025whatshouldllmsforget}
connect the recoverability axis to right-to-erasure obligations that regulated
domains, healthcare, finance, and legal, will impose whether or not the field is
ready. \citet{dmitrenko2026memorycomponents} adds a dimension the agent-memory
literature ignores entirely: the license and sustainability posture of the
open-source data-infrastructure projects that back agent memory is itself an
architectural variable for always-on durability, since a memory substrate whose
upstream project is abandoned is a durability failure waiting to happen.

The open problem for this fifth program is integrative rather than local. Each import
above lands a single governance axis: belief revision for mutability, capability
security for authority, provenance DAGs and Merkle logs for provenance,
transactional recovery for recoverability, runtime verification for validation. The
thesis demands that the six axes and five invariants hold \emph{jointly} over the
full lifecycle, yet no work in the corpus composes more than two or three axes at
once, and the imports come from disciplines with incompatible assumptions: immutable
ledgers from blockchain versus cascading deletion from referential integrity,
eventual consistency from distributed systems versus the strict precondition gates of
runtime verification. The target is a unified persistent-state governance layer in
which authority monotonicity, scope non-expansion, deletion propagation, provenance
preservation, and rollback traceability are enforced together and measured together,
with the cross-surface benchmark of the third program and the recovery harness of the
first program as its evaluation. Until that layer exists, the survey's diagnosis is
also an agenda: the works we coded accumulate and retrieve state with much more
sophistication than they govern it, and closing that gap requires importing,
composing, and measuring the governance backbone that always-on agents, as
persistent-state systems, require.
